\chapter{Crystal-field fit (menu 10)}
\label{ch:menu10}
\menulabel{10}

\section{Purpose}

Fit the crystal-field parameters $B_k^q$ from sampled state energies and their
projection coefficients -- the mean-field-cost route: cheap sampled states plus
exact linear algebra. The same menu carries the projection-basis declaration
check (A5) that decides which parameters may be compared with another source.

\section{What you need}

A JSON file with \code{point\_group}, \code{J}, \code{levels} (one energy per
sampled state) and \code{coefficients} (one row per state: the amplitudes on
$|J M\rangle$ with $M = m - J$, as numbers or \code{[re, im]} pairs). The
schema, including the optional \code{declaration} and \code{compare} blocks, is
in Chapter~\ref{ch:formats}.

\lstinputlisting[style=fbkcode, firstline=54, lastline=76,
  title={The menu-10 example session (a C3 fit with declaration and comparison)}]
  {examples/expected/m10_crystal_field.txt}

\section{What you get}

The fitted $B_k^q$ (respecting the point group's allowed set: 9 for a true
$C_3$ axis, 4 for $O_h$, 27 for no symmetry), the constant, the design-matrix
condition number, rank/conditioning notes, and -- when the JSON carries a
\code{declaration} or a \code{compare} block -- the A5 section in the report:

\lstinputlisting[style=fbkcode, firstline=21, lastline=44,
  title={The A5 declaration check in the report of the same run}]
  {examples/expected/products/cf_c3.json.fbk.md}

\section{How to read the numbers}

\begin{readbox}
\begin{itemize}
  \item \textbf{Residual} is per sampled state, not a level structure.
    Diagonalising the fitted Hamiltonian must reproduce your input levels --
    the sanity check -- while the spectrum is what a measurement or a
    multireference calculation would be compared against.
  \item \textbf{Condition number and rank} tell you whether the sampling
    supports the parameters you are asking for. Sampling eigenstates alone is
    rank-deficient: the members of a time-reversal pair share every
    expectation value, so each pair contributes one equation instead of two;
    sample oriented (non-eigen-) states, as the example does.
  \item \textbf{Comparability}: numbers computed under a different operator
    convention, projection manifold, unit or $z$-axis are not comparable. The
    A5 check refuses to compare transverse/high-order parameters across
    schemes because the same wavefunction moved $B_2^2$ from $+198$ to
    $-43$~cm$^{-1}$ between projection choices; only the leading axial
    $(2,0)$ may be compared roughly (the one measured example moved it by 3%).
\end{itemize}
\end{readbox}

\section{Worked example}

The example input is generated by
\file{examples/generate\_inputs.py}: 24 oriented states sampled from a known
$C_3$ Hamiltonian with $B_2^0 = 1000$, $B_4^0 = 10$, $B_4^3 = 50$,
$B_6^0 = 1$~cm$^{-1}$ (the states are random vectors in the $|J M\rangle$
basis, quantised so the example input is platform-independent). The fit
recovers the four values essentially exactly -- the printed parameters are
1000, 10, 50, 1 -- with a maximum residual of $4.5\times10^{-6}$~cm$^{-1}$
(the input quantisation) and condition number $4.2\times10^5$: a well-posed
problem. All other parameters come out at the $10^{-10}$ level or below, as
they must. The same run then declines to compare three of the bundled
reference parameters (the fourth, $(2,0)$, is licensed for a rough
cross-scheme comparison only).

\section{Caveats, refusals and related sections}

\begin{refusalbox}
\begin{itemize}
  \item The fit does not choose your sampling, your point group or your $J$;
    it validates the algebra and reports the diagnostics. A rank-deficient or
    ill-conditioned system is flagged, and a malformed JSON is refused with a
    next-step message.
  \item Coordinate conjugation trap: pass the amplitudes themselves, not their
    complex conjugates -- conjugating the coordinates flips the sign of every
    $q < 0$ parameter while leaving the spectrum unchanged, so the fit would
    look good with wrong sine-type parameters. The fitter warns about this in
    its notes.
  \item The A5 check decides comparability from declarations alone, never from
    value agreement: two identical-looking numbers under different schemes are
    still refused.
\end{itemize}
\end{refusalbox}

Related: \menu{2} (the point group and the parameter count), \menu{11} (a
model-level estimate of the field from the geometry), and
Chapter~\ref{ch:criteria} for the operator conventions.
